% ============================================================%
%  To preview it inside the paper temporarily, replace
%    \section{Introduction}
\label{sec:intro}

\authornote{\textbf{P1 -- clinical motivation (MIDL requires it).} Accelerated MRI
shortens scans (motion, paediatrics, throughput). Undersampling makes reconstruction
ill-posed, and learned reconstructions can hallucinate structure. Cite the false
vessels and clefts reported in the 2020 fastMRI challenge
\citep{muckley2021fastmri} and the hallucination framework of \citet{bhadra2021hallucination}.
State the stakes: any change to the DC step changes what the reconstruction is allowed
to invent.}

\authornote{\textbf{P2 -- training-free diffusion priors.} One unconditional prior serves
many masks, accelerations and coil arrays \citep{jalal2021robust,chung2022scoremri,aali2025ambient},
unlike supervised unrolled networks \citep{sriram2020varnet}. Name the sampler families:
gradient guidance (DPS, \citealp{chung2023dps}), pseudo-inverse guidance
(\pigdm, \citealp{song2023pigdm}), Krylov/CG DC (DDS, \citealp{chung2024dds}),
plug-and-play (DiffPIR, \citealp{zhu2023diffpir}) and decoupled annealing
(DAPS, \citealp{zhang2025daps}). Common denominator: an \emph{isotropic} Gaussian
approximation of $p(x_0\mid x_t)$, as shown by \citet{peng2024optimal}.}

\authornote{\textbf{P3 -- the observation (teaser, \figureref{fig:teaser}).} With a fixed,
tuned configuration, PSNR \emph{decreases} as $T$ grows from \xx{} to \xx{} for \xx{} of
\xx{} samplers. Practitioners therefore have to tune $T$ per protocol. Related
evidence: SURE-based early stopping against noise overfitting in diffusion MRI
\citep{ozturkler2023smrd}, and the finite-sample analysis of miscalibrated
intermediate posteriors \citep{burns2026finite}. Make clear that this is the first
systematic step-count study on raw multi-coil MRI, and keep the claim exactly that
narrow.}

\authornote{\textbf{P4 -- key insight.} For Cartesian MRI the forward operator is diagonal
in k-space (up to coil sensitivities). Image power is far from white across spatial
frequency, while thermal noise after prewhitening is white
\citep{roemer1990phased,kellman2005snr}. The isotropic approximation therefore
over-trusts the measurements at frequencies where the signal is weak relative to the
noise, and under-trusts them where it is strong. A k-space-diagonal posterior variance
is the natural basis here. This contrasts with the wavelet basis of
\citet{peng2024optimal} and the dense Jacobian covariances of
\citet{boys2024tmpd,rozet2024mmps}.}

\authornote{\textbf{P5 -- contributions (3 bullets, each tied to a result).}
(i) A frequency-resolved DC update (\method) that is training-free, scan-calibrated
from the ACS region, closed-form for single-coil and CG-solved for SENSE, and plugs
into five samplers (\sectionref{sec:method}).
(ii) A step-count study on raw multi-coil brain and knee data, with a band-resolved
mechanism analysis (\sectionref{sec:res-steps}).
(iii) Validation on real low-field noise (M4Raw, 0.3\,T) and inside fastMRI+ pathology
regions (\sectionref{sec:res-snr}, \sectionref{sec:res-roi}).}

\begin{figure}[t]
\floatconts{fig:teaser}
{\caption{\textbf{More steps hurt isotropic data consistency.} PSNR vs.\ number of
reverse steps $T$ on raw multi-coil fastMRI brain ($R{=}\xx$, \xx{} test volumes,
median and IQR over volumes). Solid: samplers with isotropic DC, each tuned once on the
calibration split. Dashed: the same samplers with the frequency-resolved update.}}
{\figspec{4.2cm}{Line plot, x: $T\in\{10,20,50,100,200\}$ (log scale); y: PSNR (dB).
Lines: \ikdc, \pigdm, DDS, DiffPIR (solid) and +\method{} (dashed, same colour).
Shaded IQR bands. Source: \texttt{outputs/midl27/steps/*.json}.}}
\end{figure}
\section{Related Work}
\label{sec:related}

\authornote{\textbf{Diffusion posterior sampling for MRI.} Score-based CS-MRI
\citep{jalal2021robust,chung2022scoremri}; priors trained on corrupted data
\citep{aali2025ambient}; frequency-split diffusion processes that \emph{require
training}, namely HFS-SDE \citep{cao2024hfssde} and CDiffMR \citep{huang2023cdiffmr}.
Benchmarking: InverseBench includes multi-coil fastMRI with real raw measurements and
tunes each sampler per task \citep{zheng2025inversebench}; we follow its tuning
protocol. One sentence on where we differ: we are training-free and change only the
DC step.}

\authornote{\textbf{Beyond isotropic posterior covariance.} \citet{peng2024optimal} show
that DPS and \pigdm{} share an isotropic Gaussian approximation and replace it with a
covariance that is diagonal in a wavelet basis and learned or estimated offline. TMPD
\citep{boys2024tmpd} and MMPS \citep{rozet2024mmps} use the denoiser Jacobian (MMPS is
also evaluated on fastMRI knee). CA-DPS \citep{hamidi2024cadps} estimates a diagonal
covariance online. \citet{schioppa2026covariance} use a Fourier-structured covariance
for unconditional sampling. \emph{Difference (state it precisely):} our covariance is
diagonal in the basis that diagonalises the MRI operator, is calibrated per scan from
the ACS region with no extra network evaluations, and is evaluated on raw multi-coil
and low-field data.}

\authornote{\textbf{Frequency scheduling and adaptive DC.} \citet{tian2026frequency}
expose measurement frequencies according to the diffusion noise level and return
progressively to full-band consistency, on natural images. Our frequency gate is the
k-space instance of this idea, and we compare against their schedule
(\tableref{tab:ablation}). SMRD \citep{ozturkler2023smrd} adapts a \emph{scalar} DC
weight with SURE and stops early. Measured-subspace consistency \citep{lee2026msc} is
a plug-in range/null-space correction for MRI samplers. It is complementary to ours,
because it does not reweight frequencies. Blind or noise-adaptive samplers estimate a
scalar noise level or a blur kernel
\citep{laroche2024fastdiffem,xia2026nanhmc}.}

\authornote{\textbf{Noise in MRI.} Thermal noise is white per channel and correlated
across coils; prewhitening with the coil noise covariance is standard
\citep{roemer1990phased,pruessmann1999sense,kellman2005snr}. Magnitude images follow
Rician or noncentral-$\chi$ statistics \citep{gudbjartsson1995rician,constantinides1997snr}.
Keep this to 2--3 sentences. It justifies the white-noise likelihood in
\equationref{eq:lik} and the magnitude-domain evaluation caveat.}

%  with % ============================================================%
%  To preview it inside the paper temporarily, replace
%    \section{Introduction}
\label{sec:intro}

\authornote{\textbf{P1 -- clinical motivation (MIDL requires it).} Accelerated MRI
shortens scans (motion, paediatrics, throughput). Undersampling makes reconstruction
ill-posed, and learned reconstructions can hallucinate structure. Cite the false
vessels and clefts reported in the 2020 fastMRI challenge
\citep{muckley2021fastmri} and the hallucination framework of \citet{bhadra2021hallucination}.
State the stakes: any change to the DC step changes what the reconstruction is allowed
to invent.}

\authornote{\textbf{P2 -- training-free diffusion priors.} One unconditional prior serves
many masks, accelerations and coil arrays \citep{jalal2021robust,chung2022scoremri,aali2025ambient},
unlike supervised unrolled networks \citep{sriram2020varnet}. Name the sampler families:
gradient guidance (DPS, \citealp{chung2023dps}), pseudo-inverse guidance
(\pigdm, \citealp{song2023pigdm}), Krylov/CG DC (DDS, \citealp{chung2024dds}),
plug-and-play (DiffPIR, \citealp{zhu2023diffpir}) and decoupled annealing
(DAPS, \citealp{zhang2025daps}). Common denominator: an \emph{isotropic} Gaussian
approximation of $p(x_0\mid x_t)$, as shown by \citet{peng2024optimal}.}

\authornote{\textbf{P3 -- the observation (teaser, \figureref{fig:teaser}).} With a fixed,
tuned configuration, PSNR \emph{decreases} as $T$ grows from \xx{} to \xx{} for \xx{} of
\xx{} samplers. Practitioners therefore have to tune $T$ per protocol. Related
evidence: SURE-based early stopping against noise overfitting in diffusion MRI
\citep{ozturkler2023smrd}, and the finite-sample analysis of miscalibrated
intermediate posteriors \citep{burns2026finite}. Make clear that this is the first
systematic step-count study on raw multi-coil MRI, and keep the claim exactly that
narrow.}

\authornote{\textbf{P4 -- key insight.} For Cartesian MRI the forward operator is diagonal
in k-space (up to coil sensitivities). Image power is far from white across spatial
frequency, while thermal noise after prewhitening is white
\citep{roemer1990phased,kellman2005snr}. The isotropic approximation therefore
over-trusts the measurements at frequencies where the signal is weak relative to the
noise, and under-trusts them where it is strong. A k-space-diagonal posterior variance
is the natural basis here. This contrasts with the wavelet basis of
\citet{peng2024optimal} and the dense Jacobian covariances of
\citet{boys2024tmpd,rozet2024mmps}.}

\authornote{\textbf{P5 -- contributions (3 bullets, each tied to a result).}
(i) A frequency-resolved DC update (\method) that is training-free, scan-calibrated
from the ACS region, closed-form for single-coil and CG-solved for SENSE, and plugs
into five samplers (\sectionref{sec:method}).
(ii) A step-count study on raw multi-coil brain and knee data, with a band-resolved
mechanism analysis (\sectionref{sec:res-steps}).
(iii) Validation on real low-field noise (M4Raw, 0.3\,T) and inside fastMRI+ pathology
regions (\sectionref{sec:res-snr}, \sectionref{sec:res-roi}).}

\begin{figure}[t]
\floatconts{fig:teaser}
{\caption{\textbf{More steps hurt isotropic data consistency.} PSNR vs.\ number of
reverse steps $T$ on raw multi-coil fastMRI brain ($R{=}\xx$, \xx{} test volumes,
median and IQR over volumes). Solid: samplers with isotropic DC, each tuned once on the
calibration split. Dashed: the same samplers with the frequency-resolved update.}}
{\figspec{4.2cm}{Line plot, x: $T\in\{10,20,50,100,200\}$ (log scale); y: PSNR (dB).
Lines: \ikdc, \pigdm, DDS, DiffPIR (solid) and +\method{} (dashed, same colour).
Shaded IQR bands. Source: \texttt{outputs/midl27/steps/*.json}.}}
\end{figure}
\section{Related Work}
\label{sec:related}

\authornote{\textbf{Diffusion posterior sampling for MRI.} Score-based CS-MRI
\citep{jalal2021robust,chung2022scoremri}; priors trained on corrupted data
\citep{aali2025ambient}; frequency-split diffusion processes that \emph{require
training}, namely HFS-SDE \citep{cao2024hfssde} and CDiffMR \citep{huang2023cdiffmr}.
Benchmarking: InverseBench includes multi-coil fastMRI with real raw measurements and
tunes each sampler per task \citep{zheng2025inversebench}; we follow its tuning
protocol. One sentence on where we differ: we are training-free and change only the
DC step.}

\authornote{\textbf{Beyond isotropic posterior covariance.} \citet{peng2024optimal} show
that DPS and \pigdm{} share an isotropic Gaussian approximation and replace it with a
covariance that is diagonal in a wavelet basis and learned or estimated offline. TMPD
\citep{boys2024tmpd} and MMPS \citep{rozet2024mmps} use the denoiser Jacobian (MMPS is
also evaluated on fastMRI knee). CA-DPS \citep{hamidi2024cadps} estimates a diagonal
covariance online. \citet{schioppa2026covariance} use a Fourier-structured covariance
for unconditional sampling. \emph{Difference (state it precisely):} our covariance is
diagonal in the basis that diagonalises the MRI operator, is calibrated per scan from
the ACS region with no extra network evaluations, and is evaluated on raw multi-coil
and low-field data.}

\authornote{\textbf{Frequency scheduling and adaptive DC.} \citet{tian2026frequency}
expose measurement frequencies according to the diffusion noise level and return
progressively to full-band consistency, on natural images. Our frequency gate is the
k-space instance of this idea, and we compare against their schedule
(\tableref{tab:ablation}). SMRD \citep{ozturkler2023smrd} adapts a \emph{scalar} DC
weight with SURE and stops early. Measured-subspace consistency \citep{lee2026msc} is
a plug-in range/null-space correction for MRI samplers. It is complementary to ours,
because it does not reweight frequencies. Blind or noise-adaptive samplers estimate a
scalar noise level or a blur kernel
\citep{laroche2024fastdiffem,xia2026nanhmc}.}

\authornote{\textbf{Noise in MRI.} Thermal noise is white per channel and correlated
across coils; prewhitening with the coil noise covariance is standard
\citep{roemer1990phased,pruessmann1999sense,kellman2005snr}. Magnitude images follow
Rician or noncentral-$\chi$ statistics \citep{gudbjartsson1995rician,constantinides1997snr}.
Keep this to 2--3 sentences. It justifies the white-noise likelihood in
\equationref{eq:lik} and the magnitude-domain evaluation caveat.}

%  with % ============================================================%
%  To preview it inside the paper temporarily, replace
%    \section{Introduction}
\label{sec:intro}

\authornote{\textbf{P1 -- clinical motivation (MIDL requires it).} Accelerated MRI
shortens scans (motion, paediatrics, throughput). Undersampling makes reconstruction
ill-posed, and learned reconstructions can hallucinate structure. Cite the false
vessels and clefts reported in the 2020 fastMRI challenge
\citep{muckley2021fastmri} and the hallucination framework of \citet{bhadra2021hallucination}.
State the stakes: any change to the DC step changes what the reconstruction is allowed
to invent.}

\authornote{\textbf{P2 -- training-free diffusion priors.} One unconditional prior serves
many masks, accelerations and coil arrays \citep{jalal2021robust,chung2022scoremri,aali2025ambient},
unlike supervised unrolled networks \citep{sriram2020varnet}. Name the sampler families:
gradient guidance (DPS, \citealp{chung2023dps}), pseudo-inverse guidance
(\pigdm, \citealp{song2023pigdm}), Krylov/CG DC (DDS, \citealp{chung2024dds}),
plug-and-play (DiffPIR, \citealp{zhu2023diffpir}) and decoupled annealing
(DAPS, \citealp{zhang2025daps}). Common denominator: an \emph{isotropic} Gaussian
approximation of $p(x_0\mid x_t)$, as shown by \citet{peng2024optimal}.}

\authornote{\textbf{P3 -- the observation (teaser, \figureref{fig:teaser}).} With a fixed,
tuned configuration, PSNR \emph{decreases} as $T$ grows from \xx{} to \xx{} for \xx{} of
\xx{} samplers. Practitioners therefore have to tune $T$ per protocol. Related
evidence: SURE-based early stopping against noise overfitting in diffusion MRI
\citep{ozturkler2023smrd}, and the finite-sample analysis of miscalibrated
intermediate posteriors \citep{burns2026finite}. Make clear that this is the first
systematic step-count study on raw multi-coil MRI, and keep the claim exactly that
narrow.}

\authornote{\textbf{P4 -- key insight.} For Cartesian MRI the forward operator is diagonal
in k-space (up to coil sensitivities). Image power is far from white across spatial
frequency, while thermal noise after prewhitening is white
\citep{roemer1990phased,kellman2005snr}. The isotropic approximation therefore
over-trusts the measurements at frequencies where the signal is weak relative to the
noise, and under-trusts them where it is strong. A k-space-diagonal posterior variance
is the natural basis here. This contrasts with the wavelet basis of
\citet{peng2024optimal} and the dense Jacobian covariances of
\citet{boys2024tmpd,rozet2024mmps}.}

\authornote{\textbf{P5 -- contributions (3 bullets, each tied to a result).}
(i) A frequency-resolved DC update (\method) that is training-free, scan-calibrated
from the ACS region, closed-form for single-coil and CG-solved for SENSE, and plugs
into five samplers (\sectionref{sec:method}).
(ii) A step-count study on raw multi-coil brain and knee data, with a band-resolved
mechanism analysis (\sectionref{sec:res-steps}).
(iii) Validation on real low-field noise (M4Raw, 0.3\,T) and inside fastMRI+ pathology
regions (\sectionref{sec:res-snr}, \sectionref{sec:res-roi}).}

\begin{figure}[t]
\floatconts{fig:teaser}
{\caption{\textbf{More steps hurt isotropic data consistency.} PSNR vs.\ number of
reverse steps $T$ on raw multi-coil fastMRI brain ($R{=}\xx$, \xx{} test volumes,
median and IQR over volumes). Solid: samplers with isotropic DC, each tuned once on the
calibration split. Dashed: the same samplers with the frequency-resolved update.}}
{\figspec{4.2cm}{Line plot, x: $T\in\{10,20,50,100,200\}$ (log scale); y: PSNR (dB).
Lines: \ikdc, \pigdm, DDS, DiffPIR (solid) and +\method{} (dashed, same colour).
Shaded IQR bands. Source: \texttt{outputs/midl27/steps/*.json}.}}
\end{figure}
\section{Related Work}
\label{sec:related}

\authornote{\textbf{Diffusion posterior sampling for MRI.} Score-based CS-MRI
\citep{jalal2021robust,chung2022scoremri}; priors trained on corrupted data
\citep{aali2025ambient}; frequency-split diffusion processes that \emph{require
training}, namely HFS-SDE \citep{cao2024hfssde} and CDiffMR \citep{huang2023cdiffmr}.
Benchmarking: InverseBench includes multi-coil fastMRI with real raw measurements and
tunes each sampler per task \citep{zheng2025inversebench}; we follow its tuning
protocol. One sentence on where we differ: we are training-free and change only the
DC step.}

\authornote{\textbf{Beyond isotropic posterior covariance.} \citet{peng2024optimal} show
that DPS and \pigdm{} share an isotropic Gaussian approximation and replace it with a
covariance that is diagonal in a wavelet basis and learned or estimated offline. TMPD
\citep{boys2024tmpd} and MMPS \citep{rozet2024mmps} use the denoiser Jacobian (MMPS is
also evaluated on fastMRI knee). CA-DPS \citep{hamidi2024cadps} estimates a diagonal
covariance online. \citet{schioppa2026covariance} use a Fourier-structured covariance
for unconditional sampling. \emph{Difference (state it precisely):} our covariance is
diagonal in the basis that diagonalises the MRI operator, is calibrated per scan from
the ACS region with no extra network evaluations, and is evaluated on raw multi-coil
and low-field data.}

\authornote{\textbf{Frequency scheduling and adaptive DC.} \citet{tian2026frequency}
expose measurement frequencies according to the diffusion noise level and return
progressively to full-band consistency, on natural images. Our frequency gate is the
k-space instance of this idea, and we compare against their schedule
(\tableref{tab:ablation}). SMRD \citep{ozturkler2023smrd} adapts a \emph{scalar} DC
weight with SURE and stops early. Measured-subspace consistency \citep{lee2026msc} is
a plug-in range/null-space correction for MRI samplers. It is complementary to ours,
because it does not reweight frequencies. Blind or noise-adaptive samplers estimate a
scalar noise level or a blur kernel
\citep{laroche2024fastdiffem,xia2026nanhmc}.}

\authornote{\textbf{Noise in MRI.} Thermal noise is white per channel and correlated
across coils; prewhitening with the coil noise covariance is standard
\citep{roemer1990phased,pruessmann1999sense,kellman2005snr}. Magnitude images follow
Rician or noncentral-$\chi$ statistics \citep{gudbjartsson1995rician,constantinides1997snr}.
Keep this to 2--3 sentences. It justifies the white-noise likelihood in
\equationref{eq:lik} and the magnitude-domain evaluation caveat.}

%  with % ============================================================%
%  To preview it inside the paper temporarily, replace
%    \input{sections/introduction}\input{sections/related}
%  with \input{drafts/intro_related_REFERENCE} in main.tex.
% ============================================================

\section{Introduction}
\label{sec:intro}

Magnetic resonance imaging offers unmatched soft-tissue contrast, but its
acquisition is slow. Each phase-encode line of k-space takes one repetition, so
undersampling by a factor $R$ shortens a scan almost proportionally. Shorter
scans reduce motion artefacts, make paediatric and cardiac protocols feasible,
and increase scanner throughput. The price is an ill-posed inverse problem, and
the way a reconstruction method resolves that ill-posedness determines what it
is allowed to invent. The 2020 fastMRI challenge found that top-ranked
learned reconstructions could hallucinate vessels and clefts that radiologists
judged clinically misleading \citep{muckley2021fastmri}. This makes the
mechanism by which measurements are trusted, and not only the prior, a
clinically relevant design choice \citep{bhadra2021hallucination}.

Diffusion models have become a strong, training-free route to this problem.
A single unconditional prior, trained once, can be combined with any
undersampling mask, acceleration and coil array at inference time
\citep{jalal2021robust,chung2022scoremri,aali2025ambient}. This contrasts
with supervised unrolled networks, which must be retrained for each
configuration \citep{sriram2020varnet}. Many posterior samplers have been
proposed. Gradient guidance through the denoiser (DPS,
\citealp{chung2023dps}), pseudo-inverse guidance (\pigdm,
\citealp{song2023pigdm}), conjugate-gradient data consistency on the denoised
estimate (DDS, \citealp{chung2024dds}), plug-and-play proximal steps (DiffPIR,
\citealp{zhu2023diffpir}) and decoupled annealing (DAPS,
\citealp{zhang2025daps}) differ in many details. As \citet{peng2024optimal}
showed, however, most of them share one modelling decision: the intractable
posterior $p(x_0\mid x_t)$ is approximated by a Gaussian centred on the
denoiser output with an \emph{isotropic} covariance.

We find that this decision has a visible and practically important
consequence in accelerated MRI. On raw multi-coil fastMRI brain data, samplers
with isotropic data consistency lose quality as the number of reverse steps
grows. Tuned \pigdm{} drops by \xx\,dB from its best step count to $T{=}200$,
even when its hyperparameters are re-tuned for every $T$
(\figureref{fig:teaser}). More computation should not make a posterior sampler
worse. This behaviour forces practitioners to tune the step budget per protocol.
It is consistent with earlier reports that diffusion MRI samplers overfit
measurement noise and benefit from early stopping
\citep{ozturkler2023smrd}, and with recent theory showing that miscalibrated
intermediate posteriors make samplers sensitive to when they stop
\citep{burns2026finite}.

Our explanation is specific to MRI. For Cartesian sampling, the forward
operator is diagonal in k-space up to the coil sensitivities. After
prewhitening, thermal noise is white across k-space
\citep{roemer1990phased,kellman2005snr}, but image power is anything but white:
it decays by several orders of magnitude from the centre to the periphery. An
isotropic posterior covariance ignores this decay. It therefore over-trusts
measurements at high spatial frequencies, where the signal is comparable to the
noise, and it turns data consistency at large diffusion noise into an
ill-conditioned SENSE inversion. The fix suggested by the problem's structure is
a covariance that is diagonal in the basis that diagonalises the operator. This
differs from the wavelet basis of \citet{peng2024optimal} and from the dense
Jacobian covariances of \citet{boys2024tmpd} and \citet{rozet2024mmps}, both of
which require additional network evaluations.

We make three contributions.
\begin{itemize}
  \item \textbf{A frequency-resolved data-consistency update.} \method{}
  replaces the isotropic covariance by a per-frequency posterior variance
  derived from the image power spectrum, which is estimated from the scan's own
  acquired lines. It is training-free, closed-form for single-coil data,
  solved by conjugate gradients for SENSE, and plugs into five samplers without
  extra denoiser evaluations (\sectionref{sec:method}).
  \item \textbf{A step-count study on raw multi-coil MRI.} We document the
  degradation of isotropic samplers with $T$ under per-$T$ tuning. We show that
  the frequency-resolved update removes it with a single configuration
  (\pxx\,dB at $T{=}100$, 95\% CI \ci), and we trace the effect to
  high-frequency error accumulation with a band-resolved mechanism analysis
  (\sectionref{sec:res-steps}).
  \item \textbf{Validation where it matters clinically.} We evaluate on real
  0.3\,T low-field data (M4Raw), in a controlled SNR sweep, and inside fastMRI+
  pathology regions, where the update reduces error by \xx\% without increasing
  hallucination rates (\sectionref{sec:res-snr}, \sectionref{sec:res-roi}).
\end{itemize}

\section{Related Work}
\label{sec:related}

\paragraph{Diffusion posterior sampling for MRI.}
Score-based compressed-sensing MRI was introduced with annealed Langevin
dynamics on fastMRI \citep{jalal2021robust} and extended to parallel imaging
\citep{chung2022scoremri}. \citet{aali2025ambient} train EDM priors
\citep{karras2022edm} directly on undersampled multi-coil data; we use their
fully sampled checkpoint. Several methods change the diffusion process itself
to exploit k-space structure, for example by diffusing only high-frequency
k-space \citep{cao2024hfssde} or replacing Gaussian noise with undersampling
\citep{huang2023cdiffmr}. These require training. Our update changes only the
data-consistency step of a frozen prior. InverseBench
\citep{zheng2025inversebench} benchmarks fourteen plug-and-play diffusion
samplers on multi-coil fastMRI with a per-task tuning protocol, which we
follow.

\paragraph{Beyond isotropic posterior covariances.}
\citet{peng2024optimal} show that DPS, \pigdm{} and related methods share an
isotropic Gaussian approximation of $p(x_0\mid x_t)$. They replace it with a
maximum-likelihood covariance that is learned or estimated offline and, in its
most expressive form, is diagonal in a wavelet basis; they evaluate on natural
images. Tweedie moment projection \citep{boys2024tmpd} and moment-matching
posterior sampling \citep{rozet2024mmps} use the denoiser Jacobian to form the
covariance; the latter includes accelerated knee MRI. CA-DPS
\citep{hamidi2024cadps} estimates a diagonal covariance online from score
differences, and \citet{schioppa2026covariance} propose a Fourier-structured
covariance for unconditional sampling. Our covariance differs in three ways. It
is diagonal in the basis that diagonalises the MRI operator. It is calibrated
per scan from acquired data alone, without a training set or additional network
evaluations. And it is evaluated on raw multi-coil and low-field data.

\paragraph{Frequency scheduling and adaptive data consistency.}
\citet{tian2026frequency} expose measurement frequencies progressively
according to the diffusion noise level before returning to full-band
consistency, on natural images. The optional frequency gate in our method is
the k-space instance of this idea, and we compare against their schedule
directly. SMRD \citep{ozturkler2023smrd} adapts a scalar data-consistency
weight with Stein's unbiased risk estimate and stops early to avoid fitting
noise. Measured-subspace consistency \citep{lee2026msc} is a plug-in
range/null-space correction for MRI samplers. It is complementary to ours,
because it does not reweight frequencies. Blind or noise-adaptive samplers
estimate a scalar noise level or a blur kernel
\citep{laroche2024fastdiffem,xia2026nanhmc}, whereas we model the
\emph{signal} side of the posterior per frequency.

\paragraph{Noise in MRI.}
Receiver noise is white per channel and correlated across coils. Prewhitening
with the coil noise covariance is standard \citep{roemer1990phased,pruessmann1999sense},
and reconstruction in SNR units makes noise levels comparable across scans
\citep{kellman2005snr}. Magnitude images follow Rician or noncentral-$\chi$
statistics \citep{gudbjartsson1995rician,constantinides1997snr}, which biases
magnitude-domain metrics at low SNR; we take this into account in the low-field
evaluation.
 in main.tex.
% ============================================================

\section{Introduction}
\label{sec:intro}

Magnetic resonance imaging offers unmatched soft-tissue contrast, but its
acquisition is slow. Each phase-encode line of k-space takes one repetition, so
undersampling by a factor $R$ shortens a scan almost proportionally. Shorter
scans reduce motion artefacts, make paediatric and cardiac protocols feasible,
and increase scanner throughput. The price is an ill-posed inverse problem, and
the way a reconstruction method resolves that ill-posedness determines what it
is allowed to invent. The 2020 fastMRI challenge found that top-ranked
learned reconstructions could hallucinate vessels and clefts that radiologists
judged clinically misleading \citep{muckley2021fastmri}. This makes the
mechanism by which measurements are trusted, and not only the prior, a
clinically relevant design choice \citep{bhadra2021hallucination}.

Diffusion models have become a strong, training-free route to this problem.
A single unconditional prior, trained once, can be combined with any
undersampling mask, acceleration and coil array at inference time
\citep{jalal2021robust,chung2022scoremri,aali2025ambient}. This contrasts
with supervised unrolled networks, which must be retrained for each
configuration \citep{sriram2020varnet}. Many posterior samplers have been
proposed. Gradient guidance through the denoiser (DPS,
\citealp{chung2023dps}), pseudo-inverse guidance (\pigdm,
\citealp{song2023pigdm}), conjugate-gradient data consistency on the denoised
estimate (DDS, \citealp{chung2024dds}), plug-and-play proximal steps (DiffPIR,
\citealp{zhu2023diffpir}) and decoupled annealing (DAPS,
\citealp{zhang2025daps}) differ in many details. As \citet{peng2024optimal}
showed, however, most of them share one modelling decision: the intractable
posterior $p(x_0\mid x_t)$ is approximated by a Gaussian centred on the
denoiser output with an \emph{isotropic} covariance.

We find that this decision has a visible and practically important
consequence in accelerated MRI. On raw multi-coil fastMRI brain data, samplers
with isotropic data consistency lose quality as the number of reverse steps
grows. Tuned \pigdm{} drops by \xx\,dB from its best step count to $T{=}200$,
even when its hyperparameters are re-tuned for every $T$
(\figureref{fig:teaser}). More computation should not make a posterior sampler
worse. This behaviour forces practitioners to tune the step budget per protocol.
It is consistent with earlier reports that diffusion MRI samplers overfit
measurement noise and benefit from early stopping
\citep{ozturkler2023smrd}, and with recent theory showing that miscalibrated
intermediate posteriors make samplers sensitive to when they stop
\citep{burns2026finite}.

Our explanation is specific to MRI. For Cartesian sampling, the forward
operator is diagonal in k-space up to the coil sensitivities. After
prewhitening, thermal noise is white across k-space
\citep{roemer1990phased,kellman2005snr}, but image power is anything but white:
it decays by several orders of magnitude from the centre to the periphery. An
isotropic posterior covariance ignores this decay. It therefore over-trusts
measurements at high spatial frequencies, where the signal is comparable to the
noise, and it turns data consistency at large diffusion noise into an
ill-conditioned SENSE inversion. The fix suggested by the problem's structure is
a covariance that is diagonal in the basis that diagonalises the operator. This
differs from the wavelet basis of \citet{peng2024optimal} and from the dense
Jacobian covariances of \citet{boys2024tmpd} and \citet{rozet2024mmps}, both of
which require additional network evaluations.

We make three contributions.
\begin{itemize}
  \item \textbf{A frequency-resolved data-consistency update.} \method{}
  replaces the isotropic covariance by a per-frequency posterior variance
  derived from the image power spectrum, which is estimated from the scan's own
  acquired lines. It is training-free, closed-form for single-coil data,
  solved by conjugate gradients for SENSE, and plugs into five samplers without
  extra denoiser evaluations (\sectionref{sec:method}).
  \item \textbf{A step-count study on raw multi-coil MRI.} We document the
  degradation of isotropic samplers with $T$ under per-$T$ tuning. We show that
  the frequency-resolved update removes it with a single configuration
  (\pxx\,dB at $T{=}100$, 95\% CI \ci), and we trace the effect to
  high-frequency error accumulation with a band-resolved mechanism analysis
  (\sectionref{sec:res-steps}).
  \item \textbf{Validation where it matters clinically.} We evaluate on real
  0.3\,T low-field data (M4Raw), in a controlled SNR sweep, and inside fastMRI+
  pathology regions, where the update reduces error by \xx\% without increasing
  hallucination rates (\sectionref{sec:res-snr}, \sectionref{sec:res-roi}).
\end{itemize}

\section{Related Work}
\label{sec:related}

\paragraph{Diffusion posterior sampling for MRI.}
Score-based compressed-sensing MRI was introduced with annealed Langevin
dynamics on fastMRI \citep{jalal2021robust} and extended to parallel imaging
\citep{chung2022scoremri}. \citet{aali2025ambient} train EDM priors
\citep{karras2022edm} directly on undersampled multi-coil data; we use their
fully sampled checkpoint. Several methods change the diffusion process itself
to exploit k-space structure, for example by diffusing only high-frequency
k-space \citep{cao2024hfssde} or replacing Gaussian noise with undersampling
\citep{huang2023cdiffmr}. These require training. Our update changes only the
data-consistency step of a frozen prior. InverseBench
\citep{zheng2025inversebench} benchmarks fourteen plug-and-play diffusion
samplers on multi-coil fastMRI with a per-task tuning protocol, which we
follow.

\paragraph{Beyond isotropic posterior covariances.}
\citet{peng2024optimal} show that DPS, \pigdm{} and related methods share an
isotropic Gaussian approximation of $p(x_0\mid x_t)$. They replace it with a
maximum-likelihood covariance that is learned or estimated offline and, in its
most expressive form, is diagonal in a wavelet basis; they evaluate on natural
images. Tweedie moment projection \citep{boys2024tmpd} and moment-matching
posterior sampling \citep{rozet2024mmps} use the denoiser Jacobian to form the
covariance; the latter includes accelerated knee MRI. CA-DPS
\citep{hamidi2024cadps} estimates a diagonal covariance online from score
differences, and \citet{schioppa2026covariance} propose a Fourier-structured
covariance for unconditional sampling. Our covariance differs in three ways. It
is diagonal in the basis that diagonalises the MRI operator. It is calibrated
per scan from acquired data alone, without a training set or additional network
evaluations. And it is evaluated on raw multi-coil and low-field data.

\paragraph{Frequency scheduling and adaptive data consistency.}
\citet{tian2026frequency} expose measurement frequencies progressively
according to the diffusion noise level before returning to full-band
consistency, on natural images. The optional frequency gate in our method is
the k-space instance of this idea, and we compare against their schedule
directly. SMRD \citep{ozturkler2023smrd} adapts a scalar data-consistency
weight with Stein's unbiased risk estimate and stops early to avoid fitting
noise. Measured-subspace consistency \citep{lee2026msc} is a plug-in
range/null-space correction for MRI samplers. It is complementary to ours,
because it does not reweight frequencies. Blind or noise-adaptive samplers
estimate a scalar noise level or a blur kernel
\citep{laroche2024fastdiffem,xia2026nanhmc}, whereas we model the
\emph{signal} side of the posterior per frequency.

\paragraph{Noise in MRI.}
Receiver noise is white per channel and correlated across coils. Prewhitening
with the coil noise covariance is standard \citep{roemer1990phased,pruessmann1999sense},
and reconstruction in SNR units makes noise levels comparable across scans
\citep{kellman2005snr}. Magnitude images follow Rician or noncentral-$\chi$
statistics \citep{gudbjartsson1995rician,constantinides1997snr}, which biases
magnitude-domain metrics at low SNR; we take this into account in the low-field
evaluation.
 in main.tex.
% ============================================================

\section{Introduction}
\label{sec:intro}

Magnetic resonance imaging offers unmatched soft-tissue contrast, but its
acquisition is slow. Each phase-encode line of k-space takes one repetition, so
undersampling by a factor $R$ shortens a scan almost proportionally. Shorter
scans reduce motion artefacts, make paediatric and cardiac protocols feasible,
and increase scanner throughput. The price is an ill-posed inverse problem, and
the way a reconstruction method resolves that ill-posedness determines what it
is allowed to invent. The 2020 fastMRI challenge found that top-ranked
learned reconstructions could hallucinate vessels and clefts that radiologists
judged clinically misleading \citep{muckley2021fastmri}. This makes the
mechanism by which measurements are trusted, and not only the prior, a
clinically relevant design choice \citep{bhadra2021hallucination}.

Diffusion models have become a strong, training-free route to this problem.
A single unconditional prior, trained once, can be combined with any
undersampling mask, acceleration and coil array at inference time
\citep{jalal2021robust,chung2022scoremri,aali2025ambient}. This contrasts
with supervised unrolled networks, which must be retrained for each
configuration \citep{sriram2020varnet}. Many posterior samplers have been
proposed. Gradient guidance through the denoiser (DPS,
\citealp{chung2023dps}), pseudo-inverse guidance (\pigdm,
\citealp{song2023pigdm}), conjugate-gradient data consistency on the denoised
estimate (DDS, \citealp{chung2024dds}), plug-and-play proximal steps (DiffPIR,
\citealp{zhu2023diffpir}) and decoupled annealing (DAPS,
\citealp{zhang2025daps}) differ in many details. As \citet{peng2024optimal}
showed, however, most of them share one modelling decision: the intractable
posterior $p(x_0\mid x_t)$ is approximated by a Gaussian centred on the
denoiser output with an \emph{isotropic} covariance.

We find that this decision has a visible and practically important
consequence in accelerated MRI. On raw multi-coil fastMRI brain data, samplers
with isotropic data consistency lose quality as the number of reverse steps
grows. Tuned \pigdm{} drops by \xx\,dB from its best step count to $T{=}200$,
even when its hyperparameters are re-tuned for every $T$
(\figureref{fig:teaser}). More computation should not make a posterior sampler
worse. This behaviour forces practitioners to tune the step budget per protocol.
It is consistent with earlier reports that diffusion MRI samplers overfit
measurement noise and benefit from early stopping
\citep{ozturkler2023smrd}, and with recent theory showing that miscalibrated
intermediate posteriors make samplers sensitive to when they stop
\citep{burns2026finite}.

Our explanation is specific to MRI. For Cartesian sampling, the forward
operator is diagonal in k-space up to the coil sensitivities. After
prewhitening, thermal noise is white across k-space
\citep{roemer1990phased,kellman2005snr}, but image power is anything but white:
it decays by several orders of magnitude from the centre to the periphery. An
isotropic posterior covariance ignores this decay. It therefore over-trusts
measurements at high spatial frequencies, where the signal is comparable to the
noise, and it turns data consistency at large diffusion noise into an
ill-conditioned SENSE inversion. The fix suggested by the problem's structure is
a covariance that is diagonal in the basis that diagonalises the operator. This
differs from the wavelet basis of \citet{peng2024optimal} and from the dense
Jacobian covariances of \citet{boys2024tmpd} and \citet{rozet2024mmps}, both of
which require additional network evaluations.

We make three contributions.
\begin{itemize}
  \item \textbf{A frequency-resolved data-consistency update.} \method{}
  replaces the isotropic covariance by a per-frequency posterior variance
  derived from the image power spectrum, which is estimated from the scan's own
  acquired lines. It is training-free, closed-form for single-coil data,
  solved by conjugate gradients for SENSE, and plugs into five samplers without
  extra denoiser evaluations (\sectionref{sec:method}).
  \item \textbf{A step-count study on raw multi-coil MRI.} We document the
  degradation of isotropic samplers with $T$ under per-$T$ tuning. We show that
  the frequency-resolved update removes it with a single configuration
  (\pxx\,dB at $T{=}100$, 95\% CI \ci), and we trace the effect to
  high-frequency error accumulation with a band-resolved mechanism analysis
  (\sectionref{sec:res-steps}).
  \item \textbf{Validation where it matters clinically.} We evaluate on real
  0.3\,T low-field data (M4Raw), in a controlled SNR sweep, and inside fastMRI+
  pathology regions, where the update reduces error by \xx\% without increasing
  hallucination rates (\sectionref{sec:res-snr}, \sectionref{sec:res-roi}).
\end{itemize}

\section{Related Work}
\label{sec:related}

\paragraph{Diffusion posterior sampling for MRI.}
Score-based compressed-sensing MRI was introduced with annealed Langevin
dynamics on fastMRI \citep{jalal2021robust} and extended to parallel imaging
\citep{chung2022scoremri}. \citet{aali2025ambient} train EDM priors
\citep{karras2022edm} directly on undersampled multi-coil data; we use their
fully sampled checkpoint. Several methods change the diffusion process itself
to exploit k-space structure, for example by diffusing only high-frequency
k-space \citep{cao2024hfssde} or replacing Gaussian noise with undersampling
\citep{huang2023cdiffmr}. These require training. Our update changes only the
data-consistency step of a frozen prior. InverseBench
\citep{zheng2025inversebench} benchmarks fourteen plug-and-play diffusion
samplers on multi-coil fastMRI with a per-task tuning protocol, which we
follow.

\paragraph{Beyond isotropic posterior covariances.}
\citet{peng2024optimal} show that DPS, \pigdm{} and related methods share an
isotropic Gaussian approximation of $p(x_0\mid x_t)$. They replace it with a
maximum-likelihood covariance that is learned or estimated offline and, in its
most expressive form, is diagonal in a wavelet basis; they evaluate on natural
images. Tweedie moment projection \citep{boys2024tmpd} and moment-matching
posterior sampling \citep{rozet2024mmps} use the denoiser Jacobian to form the
covariance; the latter includes accelerated knee MRI. CA-DPS
\citep{hamidi2024cadps} estimates a diagonal covariance online from score
differences, and \citet{schioppa2026covariance} propose a Fourier-structured
covariance for unconditional sampling. Our covariance differs in three ways. It
is diagonal in the basis that diagonalises the MRI operator. It is calibrated
per scan from acquired data alone, without a training set or additional network
evaluations. And it is evaluated on raw multi-coil and low-field data.

\paragraph{Frequency scheduling and adaptive data consistency.}
\citet{tian2026frequency} expose measurement frequencies progressively
according to the diffusion noise level before returning to full-band
consistency, on natural images. The optional frequency gate in our method is
the k-space instance of this idea, and we compare against their schedule
directly. SMRD \citep{ozturkler2023smrd} adapts a scalar data-consistency
weight with Stein's unbiased risk estimate and stops early to avoid fitting
noise. Measured-subspace consistency \citep{lee2026msc} is a plug-in
range/null-space correction for MRI samplers. It is complementary to ours,
because it does not reweight frequencies. Blind or noise-adaptive samplers
estimate a scalar noise level or a blur kernel
\citep{laroche2024fastdiffem,xia2026nanhmc}, whereas we model the
\emph{signal} side of the posterior per frequency.

\paragraph{Noise in MRI.}
Receiver noise is white per channel and correlated across coils. Prewhitening
with the coil noise covariance is standard \citep{roemer1990phased,pruessmann1999sense},
and reconstruction in SNR units makes noise levels comparable across scans
\citep{kellman2005snr}. Magnitude images follow Rician or noncentral-$\chi$
statistics \citep{gudbjartsson1995rician,constantinides1997snr}, which biases
magnitude-domain metrics at low SNR; we take this into account in the low-field
evaluation.
 in main.tex.
% ============================================================

\section{Introduction}
\label{sec:intro}

Magnetic resonance imaging offers unmatched soft-tissue contrast, but its
acquisition is slow. Each phase-encode line of k-space takes one repetition, so
undersampling by a factor $R$ shortens a scan almost proportionally. Shorter
scans reduce motion artefacts, make paediatric and cardiac protocols feasible,
and increase scanner throughput. The price is an ill-posed inverse problem, and
the way a reconstruction method resolves that ill-posedness determines what it
is allowed to invent. The 2020 fastMRI challenge found that top-ranked
learned reconstructions could hallucinate vessels and clefts that radiologists
judged clinically misleading \citep{muckley2021fastmri}. This makes the
mechanism by which measurements are trusted, and not only the prior, a
clinically relevant design choice \citep{bhadra2021hallucination}.

Diffusion models have become a strong, training-free route to this problem.
A single unconditional prior, trained once, can be combined with any
undersampling mask, acceleration and coil array at inference time
\citep{jalal2021robust,chung2022scoremri,aali2025ambient}. This contrasts
with supervised unrolled networks, which must be retrained for each
configuration \citep{sriram2020varnet}. Many posterior samplers have been
proposed. Gradient guidance through the denoiser (DPS,
\citealp{chung2023dps}), pseudo-inverse guidance (\pigdm,
\citealp{song2023pigdm}), conjugate-gradient data consistency on the denoised
estimate (DDS, \citealp{chung2024dds}), plug-and-play proximal steps (DiffPIR,
\citealp{zhu2023diffpir}) and decoupled annealing (DAPS,
\citealp{zhang2025daps}) differ in many details. As \citet{peng2024optimal}
showed, however, most of them share one modelling decision: the intractable
posterior $p(x_0\mid x_t)$ is approximated by a Gaussian centred on the
denoiser output with an \emph{isotropic} covariance.

We find that this decision has a visible and practically important
consequence in accelerated MRI. On raw multi-coil fastMRI brain data, samplers
with isotropic data consistency lose quality as the number of reverse steps
grows. Tuned \pigdm{} drops by \xx\,dB from its best step count to $T{=}200$,
even when its hyperparameters are re-tuned for every $T$
(\figureref{fig:teaser}). More computation should not make a posterior sampler
worse. This behaviour forces practitioners to tune the step budget per protocol.
It is consistent with earlier reports that diffusion MRI samplers overfit
measurement noise and benefit from early stopping
\citep{ozturkler2023smrd}, and with recent theory showing that miscalibrated
intermediate posteriors make samplers sensitive to when they stop
\citep{burns2026finite}.

Our explanation is specific to MRI. For Cartesian sampling, the forward
operator is diagonal in k-space up to the coil sensitivities. After
prewhitening, thermal noise is white across k-space
\citep{roemer1990phased,kellman2005snr}, but image power is anything but white:
it decays by several orders of magnitude from the centre to the periphery. An
isotropic posterior covariance ignores this decay. It therefore over-trusts
measurements at high spatial frequencies, where the signal is comparable to the
noise, and it turns data consistency at large diffusion noise into an
ill-conditioned SENSE inversion. The fix suggested by the problem's structure is
a covariance that is diagonal in the basis that diagonalises the operator. This
differs from the wavelet basis of \citet{peng2024optimal} and from the dense
Jacobian covariances of \citet{boys2024tmpd} and \citet{rozet2024mmps}, both of
which require additional network evaluations.

We make three contributions.
\begin{itemize}
  \item \textbf{A frequency-resolved data-consistency update.} \method{}
  replaces the isotropic covariance by a per-frequency posterior variance
  derived from the image power spectrum, which is estimated from the scan's own
  acquired lines. It is training-free, closed-form for single-coil data,
  solved by conjugate gradients for SENSE, and plugs into five samplers without
  extra denoiser evaluations (\sectionref{sec:method}).
  \item \textbf{A step-count study on raw multi-coil MRI.} We document the
  degradation of isotropic samplers with $T$ under per-$T$ tuning. We show that
  the frequency-resolved update removes it with a single configuration
  (\pxx\,dB at $T{=}100$, 95\% CI \ci), and we trace the effect to
  high-frequency error accumulation with a band-resolved mechanism analysis
  (\sectionref{sec:res-steps}).
  \item \textbf{Validation where it matters clinically.} We evaluate on real
  0.3\,T low-field data (M4Raw), in a controlled SNR sweep, and inside fastMRI+
  pathology regions, where the update reduces error by \xx\% without increasing
  hallucination rates (\sectionref{sec:res-snr}, \sectionref{sec:res-roi}).
\end{itemize}

\section{Related Work}
\label{sec:related}

\paragraph{Diffusion posterior sampling for MRI.}
Score-based compressed-sensing MRI was introduced with annealed Langevin
dynamics on fastMRI \citep{jalal2021robust} and extended to parallel imaging
\citep{chung2022scoremri}. \citet{aali2025ambient} train EDM priors
\citep{karras2022edm} directly on undersampled multi-coil data; we use their
fully sampled checkpoint. Several methods change the diffusion process itself
to exploit k-space structure, for example by diffusing only high-frequency
k-space \citep{cao2024hfssde} or replacing Gaussian noise with undersampling
\citep{huang2023cdiffmr}. These require training. Our update changes only the
data-consistency step of a frozen prior. InverseBench
\citep{zheng2025inversebench} benchmarks fourteen plug-and-play diffusion
samplers on multi-coil fastMRI with a per-task tuning protocol, which we
follow.

\paragraph{Beyond isotropic posterior covariances.}
\citet{peng2024optimal} show that DPS, \pigdm{} and related methods share an
isotropic Gaussian approximation of $p(x_0\mid x_t)$. They replace it with a
maximum-likelihood covariance that is learned or estimated offline and, in its
most expressive form, is diagonal in a wavelet basis; they evaluate on natural
images. Tweedie moment projection \citep{boys2024tmpd} and moment-matching
posterior sampling \citep{rozet2024mmps} use the denoiser Jacobian to form the
covariance; the latter includes accelerated knee MRI. CA-DPS
\citep{hamidi2024cadps} estimates a diagonal covariance online from score
differences, and \citet{schioppa2026covariance} propose a Fourier-structured
covariance for unconditional sampling. Our covariance differs in three ways. It
is diagonal in the basis that diagonalises the MRI operator. It is calibrated
per scan from acquired data alone, without a training set or additional network
evaluations. And it is evaluated on raw multi-coil and low-field data.

\paragraph{Frequency scheduling and adaptive data consistency.}
\citet{tian2026frequency} expose measurement frequencies progressively
according to the diffusion noise level before returning to full-band
consistency, on natural images. The optional frequency gate in our method is
the k-space instance of this idea, and we compare against their schedule
directly. SMRD \citep{ozturkler2023smrd} adapts a scalar data-consistency
weight with Stein's unbiased risk estimate and stops early to avoid fitting
noise. Measured-subspace consistency \citep{lee2026msc} is a plug-in
range/null-space correction for MRI samplers. It is complementary to ours,
because it does not reweight frequencies. Blind or noise-adaptive samplers
estimate a scalar noise level or a blur kernel
\citep{laroche2024fastdiffem,xia2026nanhmc}, whereas we model the
\emph{signal} side of the posterior per frequency.

\paragraph{Noise in MRI.}
Receiver noise is white per channel and correlated across coils. Prewhitening
with the coil noise covariance is standard \citep{roemer1990phased,pruessmann1999sense},
and reconstruction in SNR units makes noise levels comparable across scans
\citep{kellman2005snr}. Magnitude images follow Rician or noncentral-$\chi$
statistics \citep{gudbjartsson1995rician,constantinides1997snr}, which biases
magnitude-domain metrics at low SNR; we take this into account in the low-field
evaluation.
